\شروع‌مسئله [25]
%\textit{تنها دقایقی از انتشار تمرین گذشته است که م.م، س.پ، ص.ه، ع.ب و ر.ع اقدام به ارسال پاسخ‌هایشان کرده‌اند. دستیاران آموزشی که از حجم بیکاری این افراد تعجب کرده‌اند، تصمیم گرفته‌اند سوالی بدهند که حل آن ساعت‌ها طول بکشد ...}

در یک خیابان $n$ خانه با ارتفاع های متفاوت وجود دارد.
به خانه $A$ می گوییم بر خانه $B$ مشرف است اگر ارتفاع $A$ از ارتفاع $B$ بیشتر باشد.
هدف ما در این مسئله این است که برای هر خانه تعداد خانه هایی را که در سمت راستش هستند و بر آن مشرف هستند را بیابیم.

در این سوال \textbf{حتما} باید درخت \متن‌لاتین{AVL} را پیاده‌سازی کنید.

\متن‌سیاه{ورودی:}\\
ارتفاع خانه های خیابان از چپ به راست در یک خط آمده است.

\متن‌سیاه{خروجی:}\\
به ترتیب از چپ به راست چاپ می کنیم که بر هر خانه چند خانه دیگر در سمت راستش مشرف هستند.

\متن‌سیاه{محدودیت‌ها:}\\
$1\leq n\leq 4 \times 10^4$\\
که $n$ همان تعداد خانه ها است.

\متن‌سیاه{ورودی و خروجی نمونه:}\\


%% تست اول
\begin{latin}
	\begin{tabularx}{\textwidth}{|X | X|}
		\hline
		\textbf{stdin} &
		\textbf{stdout}\\
		\hline
		0 1 2 3 4
		
		&
		
		4 3 2 1 0\\
		\hline
		
	\end{tabularx}
\end{latin}


\begin{latin}
	\begin{tabularx}{\textwidth}{|X | X|}
		\hline
		\textbf{stdin} &
		\textbf{stdout}\\
		\hline
		4 2 0 1
		
		&
		
		0 0 1 0\\
		\hline
		
	\end{tabularx}
\end{latin}


\solution{
در این مسئله کافیست شروع کنیم و از راست به چپ ارتفاع ساختمان هارا در یک tree avl بریزیم.
در هر مرحله وقتی ارتفاع یک ساختمان را اضافه کردیم کافیست چک کنیم که چند عدد بزرگتر از آن در درخت موجود است.
برای این کار علاوه بر اطلاعات درخت برای هر راس تعداد رئوسی که این راس پدرش هست را هم نگه می داریم و در $O(\log n)$ می توانیم بفهمیم چند عدد بزرگتر به درخت اضافه شده است.
همین کار را تا انتها انجام می دهیم و مرتبه زمانی الگوریتم هم $O(n\log n)$می شود.

}